In the course of the paper, we define translations between several calculi of the Linear and Graded
strands, according to the overview in Figure~\ref{fig:relationship}. As for the term and type syntax
of the different calculi, we use the same interpretation brackets across different variations of
linear- and graded-base calculi respectively, making clear from the context which interpretation we are referring to.

The interpretation
$ \textcolor{GLcolor}{\llbracket}   -   \textcolor{GLcolor}{\rrbracket} $ takes graded-base terms to linear-base terms and $ \textcolor{LGcolor}{\llparenthesis} \smidge   -   \smidge \textcolor{LGcolor}{\rrparenthesis} $ takes linear-base
terms to graded-base terms. The following mnemonics have proven useful to the authors:
\begin{itemize}
  \item Given a graded term $t$, $ \textcolor{GLcolor}{\llbracket}  \SYSTEMnt{t}  \textcolor{GLcolor}{\rrbracket} $ is its translation, into a \underline{L}inear term,
    with the square \underline{L}ilac brackets mirroring the shape of the letter L.
  \item Given a linear term $t$, $ \textcolor{LGcolor}{\llparenthesis} \smidge  \SYSTEMnt{t}  \smidge \textcolor{LGcolor}{\rrparenthesis} $ is its translation into a \underline{G}raded term,
    with the rounded \underline{G}reen parentheses mirroring the shape of the letter G.
\end{itemize}
%
The same notation is used for type and context translations.

\subsubsection{Graded Base to Linear Base translation}
\label{subsec:transl-graded-base-to-linear-base}

\NewCommandCopy{\oldSYSTEMdrule}{\SYSTEMdrule}
\NewCommandCopy{\oldEquiv}{\equiv}
\NewCommandCopy{\oldTriangleq}{\triangleq}
\renewcommand{\triangleq}{& \oldTriangleq &}
\renewcommand{\equiv}{ \oldEquiv}
\renewcommand{\SYSTEMdrule}[4][]{#3 \qquad #2} % \quad\SYSTEMdrulename{#4}}}

% Our first defintion of Graded Base to Linear Base follows:
%
\begin{gather*}
{\setlength{\arraycolsep}{0.1em}
\begin{array}{lcl@{\hspace{2em}}r}
\SYSTEMdruleTranslGtoLTmvar{} \\
\SYSTEMdruleTranslGtoLTmabs{}
\SYSTEMdruleTranslGtoLTmapp{} & \text{(Term translation)} \\[0.75em]
\SYSTEMdruleTranslGtoLTybaseTy{} \\
\SYSTEMdruleTranslGtoLTyfunTy{} & \text{(Type translation)} \\[0.75em]
\SYSTEMdruleTranslGtoLGtxtempty{} \\
\SYSTEMdruleTranslGtoLGtxtgrad{} & \text{(Context translation)}
\end{array}}
\end{gather*}
\let\SYSTEMdrule\oldSYSTEMdrule
\let\equiv\oldEquiv
\let\triangleq\oldTriangleq
%
The key idea is to use the graded modality to capture the grading of assumptions and functions.
The graded function type translates to the Linear Base function type $     \multimap     $, with its left-hand side
under $ \textcolor{coeffectColor}{\square_{ \SYSTEMnt{r} } }     $.
This translation leads to the insight that Graded Base effectively
wraps a graded modality around function inputs, capturing the usage of the parameter in the body.
Since every term in Graded Base has an associated grading in the context,
every abstraction in the translation from Graded to Linear needs to do an unboxing (let) and every application needs to do a boxing (promotion).

Contexts are translated straightforwardly. The translation can never produce a linear assumption.

\begin{example} The judgment $    \SYSTEMmv{f}  :_{\textcolor{coeffectColor}{ \SYSTEMsym{(}  \SYSTEMsym{1}  \SYSTEMsym{+}  \SYSTEMsym{1}  \SYSTEMsym{)} } }   (   \SYSTEMnt{A}  \xrightarrow{\textcolor{coeffectColor}{ \SYSTEMnt{r} } }  \SYSTEMnt{A}   )    ,   \SYSTEMmv{x}  :_{\textcolor{coeffectColor}{ \SYSTEMsym{(}   \SYSTEMnt{r}  \cdot  \SYSTEMnt{r}   \SYSTEMsym{)} } }  \SYSTEMnt{A}    \vdash_{\textsc{g} }   \SYSTEMmv{f} \,  \SYSTEMsym{(}   \SYSTEMmv{f} \,  \SYSTEMmv{x}   \SYSTEMsym{)}   :  \SYSTEMnt{A} $
in Graded Base translates to $    \SYSTEMmv{f}  : \textcolor{coeffectColor}{[}    \textcolor{coeffectColor}{\square_{ \SYSTEMnt{r} } }   \textcolor{GLcolor}{\llbracket}  \SYSTEMnt{A}  \textcolor{GLcolor}{\rrbracket}    \multimap   \textcolor{GLcolor}{\llbracket}  \SYSTEMnt{A}  \textcolor{GLcolor}{\rrbracket}   {\textcolor{coeffectColor}{]_{ \SYSTEMsym{(}  \SYSTEMsym{1}  \SYSTEMsym{+}  \SYSTEMsym{1}  \SYSTEMsym{)} } } }   ,   \SYSTEMmv{x}  : \textcolor{coeffectColor}{[}   \textcolor{GLcolor}{\llbracket}  \SYSTEMnt{A}  \textcolor{GLcolor}{\rrbracket}  {\textcolor{coeffectColor}{]_{ \SYSTEMsym{(}   \SYSTEMnt{r}  \cdot  \SYSTEMnt{r}   \SYSTEMsym{)} } } }    \vdash_{\textsc{l} }   \SYSTEMmv{f} \,   \textcolor{coeffectColor}{[}   \SYSTEMmv{f} \,   \textcolor{coeffectColor}{[}  \SYSTEMmv{x}  \textcolor{coeffectColor}{]}    \textcolor{coeffectColor}{]}    :   \textcolor{GLcolor}{\llbracket}  \SYSTEMnt{A}  \textcolor{GLcolor}{\rrbracket}  $
in Linear Base. Note the right-hand sides of an application find themselves under a promotion.
\end{example}
%
\begin{remark}
We do not translate $ \SYSTEMnt{A}  \xrightarrow{\textcolor{coeffectColor}{ \SYSTEMsym{1} } }  \SYSTEMnt{B} $ to a linear function $  \textcolor{GLcolor}{\llbracket}  \SYSTEMnt{A}  \textcolor{GLcolor}{\rrbracket}   \multimap   \textcolor{GLcolor}{\llbracket}  \SYSTEMnt{B}  \textcolor{GLcolor}{\rrbracket}  $;
  this is because $ \SYSTEMnt{A}  \xrightarrow{\textcolor{coeffectColor}{ \SYSTEMsym{1} } }  \SYSTEMnt{B} $ does not necessarily mean that the function uses its argument linearly.
  Some choices of semiring mean that $1$ does not mean `used exactly once', the simplest example
  being the single-point (trivial) semiring in which $0 = 1$ and so $1$ would mean any usage.
\end{remark}
%
We now establish key soundness results. Firstly, that the translation is type
preserving, i.e., translation of terms,
contexts, and types is such that a well-typed derivation in Graded
Base is well-typed under its translation to Linear Base. This
requires various small lemmas that the translation is a
homomorphism between graded-base contexts and linear-base contexts
which we elide here. Secondly, the translation
preserves the operational semantics, and lastly it preserves the
equational theory. Appendix~\refappendix{app:proofs-grad-to-lin} provides the full proofs and supporting lemmas.
%
\begin{restatable}{thrm}{gradToLinTranslation}[Soundness of the Graded Base to Linear Base translation]
  \begin{itemize}[itemsep=0.3em]
  \item Type preservation: $ \Delta  \vdash_{\textsc{g} }  \SYSTEMnt{t}  :  \SYSTEMnt{A}  \implies   \textcolor{GLcolor}{\llbracket}  \Delta  \textcolor{GLcolor}{\rrbracket}   \vdash_{\textsc{l} }   \textcolor{GLcolor}{\llbracket}  \SYSTEMnt{t}  \textcolor{GLcolor}{\rrbracket}   :   \textcolor{GLcolor}{\llbracket}  \SYSTEMnt{A}  \textcolor{GLcolor}{\rrbracket}  $
  \item Operational correspondence: $ \SYSTEMnt{t}  \rightsquigarrow_{\textsc{g} }  \SYSTEMnt{t'}  \implies   \textcolor{GLcolor}{\llbracket}  \SYSTEMnt{t}  \textcolor{GLcolor}{\rrbracket}   \rightsquigarrow_{\textsc{l} }^\ast   \textcolor{GLcolor}{\llbracket}  \SYSTEMnt{t'}  \textcolor{GLcolor}{\rrbracket}  $
  \item Equation preservation: $ \SYSTEMnt{t}  \equiv_{\textsc{g} }  \SYSTEMnt{t'}  \implies   \textcolor{GLcolor}{\llbracket}  \SYSTEMnt{t}  \textcolor{GLcolor}{\rrbracket}   \,\equiv_{\textsc{l} }\,   \textcolor{GLcolor}{\llbracket}  \SYSTEMnt{t'}  \textcolor{GLcolor}{\rrbracket}  $
  \end{itemize}
  \label{thrm:gradToLinTranslation}
\end{restatable}
%

\subsubsection{Linear Base to Graded Base translation}
\label{subsec:lin-to-grad-cps}

A type-preserving translation from Linear Base to Graded Base is however more difficult
since we need to find some encoding of the graded modal type $ \textcolor{LGcolor}{\llparenthesis} \smidge   \textcolor{coeffectColor}{\square_{ \SYSTEMnt{r} } }  \SYSTEMnt{A}   \smidge \textcolor{LGcolor}{\rrparenthesis} $.
From the introduction rule of the graded modality (promotion), such a translation on
terms should incur a multiplication of the context by $r$ which in Graded Base is only
provided through application of a $ \SYSTEMnt{A}  \xrightarrow{\textcolor{coeffectColor}{ \SYSTEMnt{r} } }  \SYSTEMnt{B} $ function. Thus a translation must have such
a function in a negative position, suggesting a continuation-passing style (CPS)
or Church-encoding-like form.

Taking inspiration from~\citet{lindley2025encoding} on encoding products,
we extend Plotkin's CPS translation of the call-by-name $\lambda$-calculus~\cite{plotkin1975call}
with a CPS encoding of the graded modalities of Linear Base into Graded Base:
%
\renewcommand{\triangleq}{& \oldTriangleq &}
\renewcommand{\equiv}{& \oldEquiv &}
\renewcommand{\SYSTEMdrule}[4][]{#3 \qquad #2}
%
\begin{gather*}
{\setlength{\arraycolsep}{0.1em}
\begin{array}{lcl@{\hspace{2em}}r}
\SYSTEMdruleTranslLtoGCpsTyfunTy{} \\
\SYSTEMdruleTranslLtoGCpsTyboxTy{} \\
\SYSTEMdruleTranslLtoGCpsTybaseTy{} & \text{(Type translation)} \\[0.75em]
\SYSTEMdruleTranslLtoGCpsTmvar{} \\
\SYSTEMdruleTranslLtoGCpsTmabs{} \\
\SYSTEMdruleTranslLtoGCpsTmapp{} \\
\SYSTEMdruleTranslLtoGCpsTmpr{} \\
\SYSTEMdruleTranslLtoGCpsTmlet{} & \text{(Term translation)} \\[0.75em]
\SYSTEMdruleTranslLtoGGtxtempty{} \\
\SYSTEMdruleTranslLtoGGtxtlin{} \\
\SYSTEMdruleTranslLtoGGtxtgrad{} & \text{(Context translation)}
\end{array}}
\end{gather*}
%
\let\SYSTEMdrule\oldSYSTEMdrule
\let\equiv\oldEquiv
\let\triangleq\oldTriangleq
%
\noindent
The key point is to (recursively) encode $ \textcolor{coeffectColor}{\square_{ \SYSTEMnt{r} } }  \SYSTEMnt{A} $ as
 $  (   \SYSTEMnt{A}  \xrightarrow{\textcolor{coeffectColor}{ \SYSTEMnt{r} } }   \mathrm{K}    )   \xrightarrow{\textcolor{coeffectColor}{ \SYSTEMsym{1} } }   \mathrm{K}  $ so that in the case
of promotion $ \textcolor{coeffectColor}{[}  \SYSTEMnt{t}  \textcolor{coeffectColor}{]} $ we can use application to scale
the grades of $\SYSTEMnt{t}$. Note that linear assumptions get mapped
to graded assumptions with grade $1$, though this may not mean
`linear use' (see Section~\ref{subsec:linearity-preserving} for
a solution).

\begin{restatable}{thrm}{linToGradCPSTranslation}[Soundness of the Linear Base to Graded Base translation (proofs in Appendix~\refappendix{app:proofs-lin-to-grad-cps}).]
  \begin{itemize}[itemsep=0.2em]
  \item Type preservation: $ \Gamma  \vdash_{\textsc{l} }  \SYSTEMnt{t}  :  \SYSTEMnt{A}  \implies   \textcolor{LGcolor}{\llparenthesis}  \Gamma  \textcolor{LGcolor}{\rrparenthesis}   \vdash_{\textsc{g} }   \textcolor{LGcolor}{\llparenthesis} \smidge  \SYSTEMnt{t}  \smidge \textcolor{LGcolor}{\rrparenthesis}   :   \textcolor{LGcolor}{\llparenthesis} \smidge  \SYSTEMnt{A}  \smidge \textcolor{LGcolor}{\rrparenthesis}  $
  \item Operational correspondence: $ \SYSTEMnt{t}  \rightsquigarrow_{\textsc{l} }  \SYSTEMnt{t'}  \implies   \textcolor{LGcolor}{\llparenthesis} \smidge  \SYSTEMnt{t}  \smidge \textcolor{LGcolor}{\rrparenthesis}   \rightsquigarrow_{\beta\textsc{g} }^\ast   \textcolor{LGcolor}{\llparenthesis} \smidge  \SYSTEMnt{t'}  \smidge \textcolor{LGcolor}{\rrparenthesis}  $
  \item \vspace{-0.5em} Equation preservation: $\SYSTEMnt{t}  \equiv^{\beta_\to\beta\eta_\Box}_{\textsc{l} }  \SYSTEMnt{t'} \implies   \textcolor{LGcolor}{\llparenthesis} \smidge  \SYSTEMnt{t}  \smidge \textcolor{LGcolor}{\rrparenthesis}   \equiv_{\textsc{g} }   \textcolor{LGcolor}{\llparenthesis} \smidge  \SYSTEMnt{t'}  \smidge \textcolor{LGcolor}{\rrparenthesis}  $
  \end{itemize}
  \label{thrm:linToGradCPSTranslation}
\end{restatable}
%
\noindent
As we use Plotkin's CBN CPS translation, there are two notable implications:
Firstly, regarding operational correspondence, our source uses CBN but we need full
$\beta$-reduction (via the \SYSTEMRenameRuleSemGrdcongAbs{} rule) in the target to prove simulation.
Computational adequacy then follows from the $\lambda$-calculus standardisation theorem.
Secondly, regarding equation preservation: we
preserve only $\beta$-equality for functions and not $\eta$-equality---a well-known limitation of Plotkin's translation.
However, we preserve both $\beta$- and $\eta$-equality for the graded modality,
hence we label the equational theory here as $ \equiv^{\beta_\to\beta\eta_\Box}_{\textsc{l} } $.
Section~\ref{sec:poly} considers an alternate encoding in a polymorphic Graded Base setting.
